\chapter*{Tóm tắt}

Trong thương mại điện tử, giao hàng trễ thường được xem là dấu hiệu rõ nhất cho thấy khách hàng có thể không hài lòng. Tuy nhiên, thực tế cho thấy nhiều đơn hàng vẫn được giao đúng hẹn nhưng nhận đánh giá thấp. Khách hàng có thể không hài lòng vì sản phẩm không đúng kỳ vọng, hàng bị thiếu hoặc hư hỏng, đơn có nhiều người bán nên giao nhận phức tạp, hoặc trải nghiệm mua hàng không tốt dù thời gian giao vẫn nằm trong cam kết. Đề tài gọi nhóm trường hợp này là ``bất mãn ẩn'' (\emph{hidden dissatisfaction}).

Đề tài sử dụng dữ liệu từ Olist, một nền tảng thương mại điện tử có các đơn hàng thuộc nhiều nhóm sản phẩm như đồ gia dụng, thời trang, làm đẹp, điện tử và nội thất. Mục tiêu là nhận diện sớm những đơn hàng đã giao đúng hạn nhưng có nguy cơ nhận đánh giá 1--2 sao, từ đó hỗ trợ nhân viên chăm sóc khách hàng biết nên kiểm tra những trường hợp nào trước cũng như hỗ trợ thêm thông tin để cấp quản lý ra quyết định, chiến lược trong thời đại thông tin lớn.

Tóm tắt lại, đề tài của nhóm em được thực hiện theo quy trình chuẩn như sau: hệ thống xem xét thông tin đã có của đơn hàng tại thời điểm giao hàng, chẳng hạn sản phẩm được mua, người bán, giá trị đơn, phí vận chuyển, số lượng sản phẩm, số lượng người bán và quá trình xử lý đơn. Từ các thông tin này, mô hình ước lượng mức độ rủi ro của từng sản phẩm trong đơn, sau đó tổng hợp thành mức độ ưu tiên cho cả đơn hàng. Kết quả cuối cùng là một danh sách các đơn cần được xem xét trước thay vì xử lý toàn bộ đơn hàng theo cùng một thứ tự.

Ví dụ, một khách hàng đặt hai sản phẩm từ hai người bán khác nhau và đơn hàng vẫn được giao đúng hạn. Tuy nhiên, đơn có phí vận chuyển cao, nhiều bên cùng xử lý và một người bán có quá trình giao hàng thiếu ổn định. Dù chưa có đánh giá từ khách hàng, hệ thống có thể xếp đơn này vào nhóm cần theo dõi để nhân viên kiểm tra sớm hơn các đơn có rủi ro thấp.

Do dữ liệu Olist chỉ ghi nhận đánh giá ở cấp đơn hàng, đề tài không thể khẳng định sản phẩm nào là nguyên nhân trực tiếp khiến khách hàng không hài lòng. Vì vậy, kết quả của mô hình chỉ được sử dụng như một tín hiệu hỗ trợ ưu tiên, không phải kết luận về nguyên nhân hay căn cứ để tự động xử lý khách hàng.

Kết quả phát triển ban đầu của nhóm cho thấy mô hình có thể hỗ trợ xếp hạng các đơn có nguy cơ nhận đánh giá thấp tốt hơn cách lựa chọn ngẫu nhiên. Bên cạnh mô hình dự báo, đề tài xây dựng hệ thống hỗ trợ nhân viên theo dõi danh sách rủi ro, xem thông tin đơn hàng, phân công xử lý, ghi nhận quá trình liên hệ và lưu lại lịch sử thao tác. Đồng thời, cấp quản lý có thể theo dõi tình hình xử lý, thiết lập ngưỡng ưu tiên phù hợp với năng lực của đội ngũ, quản lý các tham số vận hành và xem lại lịch sử quyết định. Qua đó, hệ thống không chỉ hỗ trợ nhân viên xử lý từng trường hợp mà còn cung cấp thông tin để nhà quản lý phân bổ nguồn lực, giám sát chất lượng xử lý và đưa ra quyết định phù hợp hơn.

\par\medskip
\noindent\textbf{Từ khóa:} thương mại điện tử, sự hài lòng khách hàng, bất mãn ẩn, học máy, Olist, dự báo rủi ro.
